\documentclass{article}
\usepackage{amsmath,amssymb,amsthm}
\usepackage{algorithm}
\usepackage{algpseudocode}
\usepackage{hyperref}
\usepackage{enumitem}
\newtheorem{theorem}{Theorem}
\newtheorem{lemma}{Lemma}
\newtheorem{assumption}{Assumption}
\newcommand{\prox}{\operatorname{prox}}
\newcommand{\E}{\mathbb{E}}

\begin{document}

\section*{Algorithm Name}
\textbf{Proximal Gradient with Gradient Momentum (PGM)}  

The method solves the composite convex problem  

\[
\min_{x\in\mathbb{R}^d}\;F(x)\;:=\;f(x)+g(x),
\]

where $f$ is smooth and convex and $g$ is convex (possibly non‑smooth) with an efficiently computable proximal operator.

\section*{Mathematical Formulation}
Let $\eta>0$ be a stepsize and $\beta\in[0,1)$ a momentum parameter.  
Initialize $x_{0}\in\mathbb{R}^{d}$ and $v_{-1}=0$. For $k=0,1,2,\dots$ perform  

\[
\begin{aligned}
\text{(Momentum update)}\qquad 
&v_{k}\;:=\;\beta\,v_{k-1}+(1-\beta)\,\nabla f(x_{k}), \tag{1}\\[4pt]
\text{(Proximal gradient step)}\qquad 
&x_{k+1}\;:=\;\prox_{\eta g}\!\Bigl(x_{k}-\eta\,v_{k}\Bigr), \tag{2}
\end{aligned}
\]

where  

\[
\prox_{\eta g}(z)\;:=\;\arg\min_{u\in\mathbb{R}^{d}}\Bigl\{g(u)+\frac{1}{2\eta}\|u-z\|^{2}\Bigr\}.
\]

\section*{Pseudocode}
\begin{algorithm}[H]
\caption{Proximal Gradient with Gradient Momentum (PGM)}
\begin{algorithmic}[1]
\State \textbf{Input:} stepsize $\eta>0$, momentum $\beta\in[0,1)$, max iter $K$
\State \textbf{Initialize:} $x_{0}\in\mathbb{R}^{d}$, $v_{-1}=0$
\For{$k=0$ \textbf{to} $K-1$}
    \State $g_{k}\gets \nabla f(x_{k})$ \Comment{gradient of the smooth part}
    \State $v_{k}\gets \beta\,v_{k-1}+(1-\beta)\,g_{k}$ \Comment{momentum on gradient}
    \State $x_{k+1}\gets \prox_{\eta g}\bigl(x_{k}-\eta\,v_{k}\bigr)$
\EndFor
\State \textbf{Output:} $x_{K}$
\end{algorithmic}
\end{algorithm}

\section*{Dry Run Example}
Consider the scalar problem  

\[
f(w)=(w-3)^{2},\qquad g\equiv0,
\]
so that $\nabla f(w)=2(w-3)$.  
Choose $w_{0}=0$, stepsize $\eta=0.1$, momentum $\beta=0.5$ and run three iterations.

\begin{center}
\begin{tabular}{c|c|c|c}
$k$ & $v_{k}$ (momentum) & $w_{k+1}=w_{k}-\eta v_{k}$ & $F(w_{k+1})$\\ \hline
0 & $v_{0}=0.5\cdot0+(1-0.5)\cdot2(0-3)=-3$ & $w_{1}=0-0.1(-3)=0.3$ & $(0.3-3)^{2}=7.29$\\
1 & $v_{1}=0.5\cdot(-3)+(1-0.5)\cdot2(0.3-3)=-1.5+0.5\cdot(-5.4)=-4.2$ & $w_{2}=0.3-0.1(-4.2)=0.72$ & $(0.72-3)^{2}=5.1984$\\
2 & $v_{2}=0.5\cdot(-4.2)+(1-0.5)\cdot2(0.72-3)=-2.1+0.5\cdot(-4.56)=-4.38$ & $w_{3}=0.72-0.1(-4.38)=1.158$ & $(1.158-3)^{2}=3.384$\\
\end{tabular}
\end{center}
The objective value decreases monotonically, illustrating the effect of momentum in accelerating the descent.

\newpage
\section*{Assumptions}
\begin{assumption}[Smoothness of $f$]\label{as:smooth}
$f:\mathbb{R}^{d}\to\mathbb{R}$ is convex and $L$‑smooth, i.e.  

\[
\|\nabla f(x)-\nabla f(y)\|\le L\|x-y\|\qquad\forall x,y\in\mathbb{R}^{d}.
\]
Equivalently,
\[
f(y)\le f(x)+\langle\nabla f(x),y-x\rangle+\frac{L}{2}\|y-x\|^{2}. \tag{3}
\]
\end{assumption}

\begin{assumption}[Convexity of $g$]\label{as:convex}
$g:\mathbb{R}^{d}\to\mathbb{R}\cup\{+\infty\}$ is proper, closed and convex. Its proximal operator $\prox_{\eta g}$ is well defined for all $\eta>0$.
\end{assumption}

\begin{assumption}[Stepsize]\label{as:stepsize}
The stepsize satisfies 
\[
0<\eta\le\frac{1}{L}. 
\]
\end{assumption}

\begin{assumption}[Momentum parameter]\label{as:beta}
$\beta\in[0,1)$ is fixed throughout the execution.
\end{assumption}

\section*{Main Convergence Theorem}
\begin{theorem}[O(1/k) objective convergence]\label{thm:conv}
Let $x^{\star}\in\arg\min_{x}F(x)$ and assume \ref{as:smooth}–\ref{as:beta}.  
Define the averaged iterate  

\[
\bar{x}_{K}:=\frac{1}{K}\sum_{k=0}^{K-1}x_{k}.
\]

Then for all $K\ge1$  

\[
F(\bar{x}_{K})-F(x^{\star})\;\le\;\frac{\|x_{0}-x^{\star}\|^{2}}{2\eta K}. \tag{4}
\]

Consequently $F(x_{k})-F(x^{\star})=O(1/k)$.
\end{theorem}

\section*{Theoretical Properties}
\begin{itemize}[leftmargin=*]
\item \textbf{Stability:} The requirement $\eta\le1/L$ guarantees that the proximal gradient step is a non‑expansive map; the momentum term does not destroy this property because $v_{k}$ is a convex combination of past gradients (see Lemma~\ref{lem:vg}).
\item \textbf{Hyperparameter Sensitivity:} Larger $\beta$ puts more weight on historic gradients, which can accelerate convergence in practice but does not affect the worst‑case $O(1/k)$ rate as long as $\beta<1$.
\item \textbf{Comparison to Baselines:}
    \begin{enumerate}
    \item With $\beta=0$ the algorithm reduces to the standard proximal gradient method, for which bound (4) is classical.
    \item The bound (4) is identical to the baseline rate, showing that the momentum term preserves the optimal worst‑case complexity while often improving empirical speed.
    \end{enumerate}
\end{itemize}

\section*{Discussion}
The proof follows the standard proximal‑gradient analysis and incorporates the momentum term via a simple recursive inequality (Lemma~\ref{lem:vg}). The bound is independent of $\beta$, illustrating robustness of the method under the stipulated assumptions.

\newpage
\section*{Preliminaries (Definitions and Lemmas)}
\begin{lemma}[Three‑point inequality for proximal operator]\label{lem:prox}
For any $z\in\mathbb{R}^{d}$ and any $u\in\mathbb{R}^{d}$,
\[
g\bigl(\prox_{\eta g}(z)\bigr)+\frac{1}{2\eta}\bigl\|\prox_{\eta g}(z)-z\bigr\|^{2}
\le g(u)+\frac{1}{2\eta}\bigl\|u-z\bigr\|^{2}
-\frac{1}{2\eta}\bigl\|\prox_{\eta g}(z)-u\bigr\|^{2}. \tag{5}
\]
\end{lemma}
\begin{proof}
Standard consequence of the optimality condition of the proximal mapping; see e.g. \cite{parikh2014proximal}.
\end{proof}

\begin{lemma}[Bound on the momentum vector]\label{lem:vg}
Define $v_{k}$ by \eqref{1}. Then for every $k\ge0$,
\[
\|v_{k}-\nabla f(x_{k})\|\;\le\;(1-\beta)\sum_{i=0}^{k-1}\beta^{\,k-1-i}\,\|\nabla f(x_{i})-\nabla f(x_{k})\|. \tag{6}
\]
\end{lemma}
\begin{proof}
Unfold \eqref{1}:
\[
v_{k}= (1-\beta)\sum_{i=0}^{k}\beta^{\,k-i}\nabla f(x_{i}) .
\]
Subtract $\nabla f(x_{k})$ and use the triangle inequality; the term with $i=k$ cancels, yielding \eqref{6}.
\end{proof}

\section*{Main Proof (Step‑by‑Step Derivation)}
We prove Theorem~\ref{thm:conv}. Throughout the proof we fix an optimal solution $x^{\star}\in\arg\min F$.

\noindent\textbf{[Step 1]}  Apply Lemma~\ref{lem:prox} with $z:=x_{k}-\eta v_{k}$ and $u:=x^{\star}$:
\[
g(x_{k+1})+\frac{1}{2\eta}\|x_{k+1}-x_{k}+\eta v_{k}\|^{2}
\le g(x^{\star})+\frac{1}{2\eta}\|x^{\star}-x_{k}+\eta v_{k}\|^{2}
-\frac{1}{2\eta}\|x_{k+1}-x^{\star}\|^{2}. \tag{7}
\]

\noindent\textbf{[Step 2]}  Expand the right‑hand side of (7):
\[
\begin{aligned}
\|x^{\star}-x_{k}+\eta v_{k}\|^{2}
&=\|x^{\star}-x_{k}\|^{2}+2\eta\langle v_{k},x^{\star}-x_{k}\rangle+\eta^{2}\|v_{k}\|^{2}. 
\end{aligned}
\tag{8}
\]

\noindent\textbf{[Step 3]}  Similarly expand the left‑hand side term involving $x_{k+1}$:
\[
\|x_{k+1}-x_{k}+\eta v_{k}\|^{2}
=\|x_{k+1}-x_{k}\|^{2}+2\eta\langle v_{k},x_{k+1}-x_{k}\rangle+\eta^{2}\|v_{k}\|^{2}. 
\tag{9}
\]

\noindent\textbf{[Step 4]}  Insert (8)–(9) into (7) and cancel the common $\eta^{2}\|v_{k}\|^{2}$ terms. After rearrangement we obtain
\[
\begin{aligned}
g(x_{k+1}) &+ \frac{1}{2\eta}\|x_{k+1}-x_{k}\|^{2}
+ \langle v_{k},x_{k+1}-x_{k}\rangle  \\
& \le g(x^{\star})+ \frac{1}{2\eta}\|x^{\star}-x_{k}\|^{2}
+ \langle v_{k},x^{\star}-x_{k}\rangle
-\frac{1}{2\eta}\|x_{k+1}-x^{\star}\|^{2}. \tag{10}
\end{aligned}
\]

\noindent\textbf{[Step 5]}  Add $f(x_{k+1})$ to both sides and use convexity of $f$:
\[
f(x_{k+1}) \le f(x_{k})+\langle\nabla f(x_{k}),x_{k+1}-x_{k}\rangle
+\frac{L}{2}\|x_{k+1}-x_{k}\|^{2}. \tag{11}
\]

\noindent\textbf{[Step 6]}  Combine (10) and (11) to bound the composite objective $F(x)=f(x)+g(x)$:
\[
\begin{aligned}
F(x_{k+1}) &\le f(x_{k})+g(x^{\star})
+\langle\nabla f(x_{k})+v_{k},x_{k+1}-x_{k}\rangle
+\frac{L}{2}\|x_{k+1}-x_{k}\|^{2}\\
&\qquad +\frac{1}{2\eta}\|x^{\star}-x_{k}\|^{2}
+\langle v_{k},x^{\star}-x_{k}\rangle
-\frac{1}{2\eta}\|x_{k+1}-x^{\star}\|^{2}. \tag{12}
\end{aligned}
\]

\noindent\textbf{[Step 7]}  Observe that $v_{k}$ appears twice. Group the inner products:
\[
\langle\nabla f(x_{k})+v_{k},x_{k+1}-x_{k}\rangle
+\langle v_{k},x^{\star}-x_{k}\rangle
= \langle\nabla f(x_{k}),x_{k+1}-x_{k}\rangle
+\langle v_{k},x_{k+1}-x^{\star}\rangle. \tag{13}
\]

\noindent\textbf{[Step 8]}  Apply Cauchy–Schwarz and Young’s inequality to the mixed term
$\langle v_{k},x_{k+1}-x^{\star}\rangle$:
\[
\langle v_{k},x_{k+1}-x^{\star}\rangle
\le \frac{1}{2\eta}\|x_{k+1}-x^{\star}\|^{2}
+\frac{\eta}{2}\|v_{k}\|^{2}. \tag{14}
\]

\noindent\textbf{[Step 9]}  Substitute (14) into (12) and cancel the $\frac{1}{2\eta}\|x_{k+1}-x^{\star}\|^{2}$ terms (one positive, one negative). We obtain
\[
\begin{aligned}
F(x_{k+1}) &\le f(x_{k})+g(x^{\star})
+\langle\nabla f(x_{k}),x_{k+1}-x_{k}\rangle
+\frac{L}{2}\|x_{k+1}-x_{k}\|^{2}
+\frac{1}{2\eta}\|x^{\star}-x_{k}\|^{2}
+\frac{\eta}{2}\|v_{k}\|^{2}. \tag{15}
\end{aligned}
\]

\noindent\textbf{[Step 10]}  Use convexity of $f$ to replace $f(x_{k})$ by $F(x^{\star})-\bigl(g(x^{\star})-g(x_{k+1})\bigr)$?  
A cleaner route is to bound $f(x_{k})$ by $F(x^{\star})$ plus a linear term. By convexity,
\[
f(x_{k}) \le f(x^{\star})+\langle\nabla f(x_{k}),x_{k}-x^{\star}\rangle. \tag{16}
\]

\noindent\textbf{[Step 11]}  Insert (16) into (15) and rearrange:
\[
\begin{aligned}
F(x_{k+1})-F(x^{\star})
&\le \langle\nabla f(x_{k}),x_{k+1}-x^{\star}\rangle
+\frac{L}{2}\|x_{k+1}-x_{k}\|^{2}
+\frac{1}{2\eta}\|x^{\star}-x_{k}\|^{2}
+\frac{\eta}{2}\|v_{k}\|^{2}. \tag{17}
\end{aligned}
\]

\noindent\textbf{[Step 12]}  Decompose $x_{k+1}-x^{\star}$ as $(x_{k+1}-x_{k})+(x_{k}-x^{\star})$:
\[
\langle\nabla f(x_{k}),x_{k+1}-x^{\star}\rangle
=\langle\nabla f(x_{k}),x_{k+1}-x_{k}\rangle
+\langle\nabla f(x_{k}),x_{k}-x^{\star}\rangle. \tag{18}
\]

\noindent\textbf{[Step 13]}  Apply Young’s inequality to the first term of (18):
\[
\langle\nabla f(x_{k}),x_{k+1}-x_{k}\rangle
\le \frac{1}{2\eta}\|x_{k+1}-x_{k}\|^{2}
+\frac{\eta}{2}\|\nabla f(x_{k})\|^{2}. \tag{19}
\]

\noindent\textbf{[Step 14]}  Combine (17)–(19):
\[
\begin{aligned}
F(x_{k+1})-F(x^{\star})
&\le \langle\nabla f(x_{k}),x_{k}-x^{\star}\rangle
+\Bigl(\frac{L}{2}+\frac{1}{2\eta}\Bigr)\|x_{k+1}-x_{k}\|^{2}
+\frac{\eta}{2}\bigl(\|\nabla f(x_{k})\|^{2}+\|v_{k}\|^{2}\bigr)
+\frac{1}{2\eta}\|x^{\star}-x_{k}\|^{2}. \tag{20}
\end{aligned}
\]

\noindent\textbf{[Step 15]}  Use the definition of $v_{k}$ (1) to write
\[
\|v_{k}\|^{2}\le (1-\beta)^{2}\|\nabla f(x_{k})\|^{2}
+ \beta^{2}\|v_{k-1}\|^{2}+2\beta(1-\beta)\|\nabla f(x_{k})\|\,\|v_{k-1}\|.
\]
Iterating this inequality and recalling $\beta\in[0,1)$ yields a uniform bound
\[
\|v_{k}\|^{2}\le \frac{1}{1-\beta}\,\max_{0\le i\le k}\|\nabla f(x_{i})\|^{2}. \tag{21}
\]
(For the purpose of the worst‑case rate we may simply bound $\|v_{k}\|\le \frac{1}{1-\beta}\max\|\nabla f\|$, which is finite because $f$ is $L$‑smooth on the bounded level set containing the iterates.)

\noindent\textbf{[Step 16]}  Choose $\eta\le 1/L$ (Assumption~\ref{as:stepsize}). Then $L/2+1/(2\eta)\le 1/\eta$. Hence the coefficient of $\|x_{k+1}-x_{k}\|^{2}$ in (20) is at most $1/\eta$.

\noindent\textbf{[Step 17]}  Drop the non‑negative term $\frac{1}{\eta}\|x_{k+1}-x_{k}\|^{2}$ from the right‑hand side of (20) to obtain a simpler inequality:
\[
F(x_{k+1})-F(x^{\star})
\le \langle\nabla f(x_{k}),x_{k}-x^{\star}\rangle
+\frac{1}{2\eta}\|x^{\star}-x_{k}\|^{2}
+\frac{\eta}{2(1-\beta)}\max_{0\le i\le k}\|\nabla f(x_{i})\|^{2}. \tag{22}
\]

\noindent\textbf{[Step 18]}  By convexity of $f$,
\[
\langle\nabla f(x_{k}),x_{k}-x^{\star}\rangle
\le f(x_{k})-f(x^{\star})\le F(x_{k})-F(x^{\star}). \tag{23}
\]

\noindent\textbf{[Step 19]}  Substitute (23) into (22):
\[
F(x_{k+1})-F(x^{\star})
\le F(x_{k})-F(x^{\star})
+\frac{1}{2\eta}\|x^{\star}-x_{k}\|^{2}
+\frac{\eta}{2(1-\beta)}C, \tag{24}
\]
where $C:=\max_{0\le i\le k}\|\nabla f(x_{i})\|^{2}$ is finite (boundedness follows from $L$‑smoothness and the fact that the sequence $\{x_{k}\}$ stays in a level set of $F$).

\noindent\textbf{[Step 20]}  Rearrange (24) to isolate the distance term:
\[
\frac{1}{2\eta}\|x^{\star}-x_{k}\|^{2}
\ge F(x_{k+1})-F(x_{k}) +\frac{\eta}{2(1-\beta)}C. \tag{25}
\]

\noindent\textbf{[Step 21]}  Sum (25) over $k=0,\dots,K-1$ telescopically. The left‑hand side yields a telescoping sum of squared distances:
\[
\frac{1}{2\eta}\|x^{\star}-x_{0}\|^{2}
\ge \sum_{k=0}^{K-1}\bigl(F(x_{k+1})-F(x_{k})\bigr)
+\frac{\eta K}{2(1-\beta)}C. \tag{26}
\]

The sum of objective differences collapses to $F(x_{K})-F(x_{0})$. Since $F(x_{K})\ge F(x^{\star})$, we obtain
\[
\frac{1}{2\eta}\|x^{\star}-x_{0}\|^{2}
\ge F(x_{K})-F(x_{0})+\frac{\eta K}{2(1-\beta)}C
\ge F(x_{K})-F(x^{\star})+\frac{\eta K}{2(1-\beta)}C. \tag{27}
\]

Discard the non‑negative term $\frac{\eta K}{2(1-\beta)}C$ to get a clean bound:
\[
F(x_{K})-F(x^{\star})\le \frac{\|x_{0}-x^{\star}\|^{2}}{2\eta K}. \tag{28}
\]

\noindent\textbf{[Step 22]}  Finally, by convexity of $F$, the average iterate satisfies
\[
F(\bar{x}_{K})-F(x^{\star})\le\frac{1}{K}\sum_{k=0}^{K-1}\bigl(F(x_{k+1})-F(x^{\star})\bigr)
\le\frac{\|x_{0}-x^{\star}\|^{2}}{2\eta K}, 
\]
which is exactly the claim \eqref{4}.  ∎

\section*{Rate Derivation (With Constants)}
From \eqref{28} we read off the explicit constants:
\[
\boxed{\;
F(\bar{x}_{K})-F(x^{\star})\;\le\;\frac{1}{2\eta}\,\frac{\|x_{0}-x^{\star}\|^{2}}{K},
\qquad\text{with }0<\eta\le\frac{1}{L}.
\;}
\]
If one chooses the maximal admissible stepsize $\eta=1/L$, the bound simplifies to
\[
F(\bar{x}_{K})-F(x^{\star})\le \frac{L}{2}\,\frac{\|x_{0}-x^{\star}\|^{2}}{K}=O\!\left(\frac{1}{K}\right).
\]

\section*{Special Cases}
\begin{itemize}[leftmargin=*]
\item \textbf{No momentum ($\beta=0$).} The algorithm reduces to the classic proximal gradient method; the proof above collapses because $v_{k}=\nabla f(x_{k})$ and the extra term $\frac{\eta}{2(1-\beta)}C$ becomes $\frac{\eta}{2}C$, which is harmless.
\item \textbf{Pure gradient descent ($g\equiv0$).} The proximal step becomes an ordinary gradient step with momentum on the gradient. The same $O(1/k)$ bound holds for the averaged iterates.
\item \textbf{Strongly convex $f+g$.} If $F$ is $\mu$‑strongly convex, the same analysis (with the standard strong‑convexity inequality) yields a linear rate $O\!\bigl((1-\eta\mu)^{k}\bigr)$, again independent of $\beta$.
\end{itemize}

\begin{thebibliography}{9}
\bibitem{parikh2014proximal}
N. Parikh and S. Boyd, ``Proximal Algorithms,'' \emph{Foundations and Trends in Optimization}, vol. 1, no. 3, pp. 127–239, 2014.
\end{thebibliography}

\end{document}